Nullniveau der Lageenergie: der tiefste Punkt des Tals. Zustand 1: auf der Kuppe (Lageenergie und kinetische Energie, beide zählen). Zustand 2: im Tal (nur kinetische Energie). Zustand 3: wo sie anhält (nur Lageenergie).

\begin{abcliste}
\abc
Sei $v_0$ ihre Geschwindigkeit auf der Kuppe, $h_1$ deren Höhe über dem Tal und $v_1$ ihre Geschwindigkeit im Tal. Zwischen den Zuständen 1 und 2:
\begin{align}
 m\,g\,h_1 + \frac{1}{2}\,m\,v_0^2 &= \frac{1}{2}\,m\,v_1^2 \\
 v_1 &= \vaF \\
   &= \sqrt{\left(\vK\right)^2+2\cdot\g\cdot\hT} \\
   &= \va \approx \result{\vaP{0}{2}}
\end{align}
Geschwindigkeiten addieren sich nicht, Energien schon.

\abc
Sei $h_2$ die Höhe, wo sie anhält. Zwischen den Zuständen 1 und 3 wird die ganze Energie zu Lageenergie:
\begin{align}
 m\,g\,h_1 + \frac{1}{2}\,m\,v_0^2 &= m\,g\,h_2 \\
 h_2 &= \hzF \\
   &= \hT+\frac{\left(\vK\right)^2}{2\cdot\g} \\
   &= \hz \approx \result{\hzP{0}{3}}
\end{align}
Höher als dort, wo sie gestartet ist: Ihre kinetische Energie am Anfang trägt sie weiter hinauf.
\end{abcliste}
